% =============================================================================
% Chapter 12: Benchmarking — ML Systems Lecture Slides
% =============================================================================
% IMAGES NEEDED (SVG → PDF):
%   - three-dimensional-framework.pdf   - benchmark-anatomy-pipeline.pdf
%   - mlperf-scenarios-quadrant.pdf     - latency-distribution-percentiles.pdf
%   - latency-breakdown-waterfall.pdf   - measurement-methodology.pdf
%   - benchmark-gaming-antipatterns.pdf - benchmark-granularity.pdf
%   - benchmark-evolution-timeline.pdf
% =============================================================================
\documentclass[aspectratio=169, 12pt]{beamer}
\usepackage{../../assets/beamerthememlsys}

\mlsyssetup{
  volume   = {Volume I},
  chapter  = {Chapter 12},
  logo = {../../assets/img/logo-mlsysbook.png},
  instlogo = {../../assets/img/logo-harvard.png},
  chaptertitle = {Benchmarking},
}

% --- Fonts ---
\usepackage{fontspec}
\setsansfont{Helvetica Neue}[
  BoldFont={Helvetica Neue Bold},
  ItalicFont={Helvetica Neue Italic},
  BoldItalicFont={Helvetica Neue Bold Italic},
]
\setmonofont{JetBrains Mono}[Scale=0.85]

% --- Packages ---
\usepackage{booktabs}
\usepackage{amsmath}

% --- Image paths ---
\graphicspath{
  {images/}
}

% --- Chapter-specific macros ---
\newcommand{\DAM}{D\raisebox{0.08em}{\tiny$\bullet$}A\raisebox{0.08em}{\tiny$\bullet$}M}

% --- Helper: safe image include ---
\newcommand{\safeimg}[2][width=\textwidth,keepaspectratio]{%
  \IfFileExists{images/#2}{\includegraphics[#1]{#2}}{%
    \IfFileExists{#2}{\includegraphics[#1]{#2}}{%
      \fbox{\parbox[c][2.5cm][c]{0.85\linewidth}{\centering\footnotesize\textcolor{midgray}{[Missing image]}}}%
    }%
  }%
}

% --- Section count for navigation (must match actual \section{} count) ---
\setcounter{mlsystotalsections}{9}

\title{Benchmarking}
\author{Vijay Janapa Reddi}
\institute{Harvard University}
\date{}

\begin{document}

% =============================================================================
% TITLE SLIDE
% =============================================================================
\mlsystitle{Benchmarking}{Measuring What Matters}{cover_ai_benchmarking.png}

% =============================================================================
% LEARNING OBJECTIVES
% =============================================================================
\begin{frame}{Learning Objectives}
\note{
% -- LINK: Ch10 compressed, Ch11 accelerated; this chapter validates the gains
Students optimized models and hardware for three chapters. This slide frames
the validation question: did those optimizations actually work?

% -- NARRATE: Walk through objectives
Read each objective. ``This chapter converts theoretical claims into verified
engineering knowledge. The 2--10x gap between lab and production is real.''

% -- ENGAGE: Surface personal experience
Ask: ``How do you know your optimizations actually worked?''

% -- WARN: Students think benchmarking is just running a test
Common error: ``I measured it once, it was fast.'' Correct framing:
rigorous benchmarking requires statistical methods, environment control,
and multi-dimensional evaluation.

% -- FLEX: [CORE] Objectives set the chapter roadmap
[CORE] This slide frames every subsequent concept.
IF AHEAD: ``Which objective surprises you?''
IF SHORT: Read objectives quickly.
}

\footnotesize
\begin{enumerate}\setlength\itemsep{0pt}
  \item Explain the \textbf{three-dimensional benchmarking framework} (system, model, data)
  \item Compare \textbf{training} and \textbf{inference} benchmarking through their distinct metrics
  \item Select appropriate \textbf{granularity levels} (micro, macro, end-to-end)
  \item Apply \textbf{MLPerf} standards across training, inference, and power
  \item Design benchmark protocols with \textbf{statistical rigor} (CI, warm-up, environment control)
  \item Distinguish \textbf{laboratory results} from production requirements
  \item Critique benchmark limitations including \textbf{Goodhart's Law} and \textbf{gaming}
\end{enumerate}

\end{frame}

\begin{frame}{Visual Language}
\note{
% -- NARRATE: Brief color code orientation
In this chapter, red highlights the benchmark-production gap, blue highlights
measurement methodology, green highlights validated results.

% -- FLEX: [OPTIONAL] Skip if students know the color system
[OPTIONAL] Say ``same legend'' and advance.
}

\small
Throughout this course, colors carry meaning:

\vspace{0.3cm}
\begin{columns}[T]
  \begin{column}{0.45\textwidth}
    \begin{mlsyscard}{computestroke}
    \textbf{Blue} --- Compute / Processing\\
    {\footnotesize GPU ops, forward/backward pass, inference}
    \end{mlsyscard}
    \vspace{0.15cm}
    \begin{mlsyscard}{datastroke}
    \textbf{Green} --- Data / Memory\\
    {\footnotesize Data flow, caches, healthy paths}
    \end{mlsyscard}
  \end{column}
  \begin{column}{0.45\textwidth}
    \begin{mlsyscard}{routingstroke}
    \textbf{Orange} --- Routing / Scheduling\\
    {\footnotesize Load balancers, batch windows}
    \end{mlsyscard}
    \vspace{0.15cm}
    \begin{mlsyscard}{errorstroke}
    \textbf{Red} --- Error / Cost / Bottleneck\\
    {\footnotesize Loss, decode phase, waste}
    \end{mlsyscard}
  \end{column}
\end{columns}
\end{frame}


% =============================================================================
\section{Why Benchmark?}
% =============================================================================

\begin{frame}{The Benchmark--Production Gap}
\note{
% -- LINK: Learning objectives promised validation; this shows why it is needed
The objectives stated benchmarking converts claims to knowledge. This slide
reveals the motivating problem: a 2--10x gap between lab and production.

% -- NARRATE: Walk through the two-column layout
Left: ``Data selection, compression, hardware --- all promised improvements.''
Red card: ``The gap is routinely 2--10x.'' Right table: ``Lab: 15 ms, 92\%,
800 QPS. Production: 150--200 ms, 78--82\%, 400--500 QPS.'' Why: ``Concurrent
load, data drift, preprocessing overhead, thermal throttling.''

% -- ENGAGE: Surface personal experience with the gap
Ask: ``Has anyone run a model that worked great in testing but failed in
production?'' Cold-call 2 students.

% -- WARN: Students trust single-environment benchmarks
Common error: ``It worked on my laptop.'' Correct framing: controlled
benchmarks systematically omit production realities.

% -- FLEX: [CORE] The gap motivates the entire chapter
[CORE] Essential --- every technique addresses this gap.
IF AHEAD: ``What is the biggest source of the gap?''
IF SHORT: Show the table, state the 2--10x number.
}

\footnotesize
\begin{columns}[T]
  \begin{column}{0.52\textwidth}
    Preceding chapters promised improvements:
    \begin{itemize}\setlength\itemsep{0pt}
      \item \textbf{Data selection}: more efficient training
      \item \textbf{Model compression}: smaller, faster models
      \item \textbf{Hardware acceleration}: higher throughput
    \end{itemize}

    \vspace{0.1cm}
    \begin{mlsyscard}{errorstroke}
    \textbf{The reality}: The gap is routinely \textbf{2$\times$--10$\times$}.\\[0.05cm]
    A quantized model may benchmark 2$\times$ faster but show \emph{no improvement} under real traffic.
    \end{mlsyscard}

    \vspace{0.05cm}
    \mlsysconcept{Core question:}{Did optimizations deliver under \emph{realistic} conditions?}
  \end{column}
  \begin{column}{0.45\textwidth}
    \renewcommand{\arraystretch}{1.15}
    \begin{tabular}{@{}lll@{}}
      \toprule
      \textbf{Claim} & \textbf{Lab} & \textbf{Prod.} \\
      \midrule
      Latency & 15 ms & 150--200 ms \\
      Accuracy & 92\% & 78--82\% \\
      Throughput & 800 QPS & 400--500 QPS \\
      \bottomrule
    \end{tabular}

    \vspace{0.15cm}
    \textcolor{errorstroke}{\textbf{Why?}} Concurrent load, data drift, preprocessing overhead, thermal throttling.
  \end{column}
\end{columns}

\end{frame}

\begin{frame}{Three Principles of Rigorous Benchmarking}
\note{
% -- LINK: The gap showed the problem; these principles show how to avoid it
The 2--10x gap was the problem. These three principles separate rigorous
engineering from guesswork.

% -- NARRATE: Walk through each principle
``First: benchmarks are proxies, not truth --- 10,000 samples/s offline may
mean only 200 QPS with latency constraints.'' ``Second: Goodhart's Law ---
when a measure becomes a target, it ceases to be good.'' ``Third: end-to-end
beats component metrics --- 3x inference speedup yields only 1.3x system gain.''

% -- ENGAGE: Can students give an example of Goodhart's Law in ML?
Ask: ``Give an example of Goodhart's Law in ML.'' Expected: optimizing for
ImageNet accuracy leads to models that fail on real-world images.

% -- WARN: Students optimize for the metric, not the goal
Common error: climbing leaderboards while degrading real-world performance.
Correct framing: the benchmark is a proxy. Always validate end-to-end.

% -- FLEX: [CORE] These principles are referenced throughout the chapter
[CORE] Students apply these when critiquing benchmarks later.
IF AHEAD: ``How does MLPerf resist Goodhart's Law?''
IF SHORT: State all three principles without elaboration.
}

\footnotesize
\textbf{1.\ Benchmarks are proxies, not truth.}\\
A system achieving 10,000 samples/s in Offline mode might achieve only 200 QPS with latency constraints. The question is always what the benchmark does \emph{not} measure.

\vspace{0.1cm}
\textbf{2.\ Goodhart's Law applies everywhere.}\\
``When a measure becomes a target, it ceases to be a good measure.'' Teams optimizing for benchmark rankings produce systems that excel in evaluation but fail in production.

\vspace{0.1cm}
\textbf{3.\ End-to-end beats component metrics.}\\
Vendors report 5--10 ms model inference. Production latency is 50--100 ms total. A 3$\times$ inference speedup yields only 1.3$\times$ end-to-end improvement.

\vspace{0.1cm}
\begin{mlsyscard}{crimson}
\textbf{Benchmarking} converts theoretical claims into verified engineering knowledge.
\end{mlsyscard}

\end{frame}

% =============================================================================
\section{Historical Foundations}
% =============================================================================

\begin{frame}{From Whetstone to MLPerf}
\note{
% -- LINK: Principles were abstract; history shows them in action
The three principles were theoretical. This timeline shows how each
generation of benchmarks learned from the gaming of its predecessor.

% -- NARRATE: Walk through the four eras on the timeline
``Whetstone (1970s): floating-point loops became a compiler optimization test.
Dhrystone (1984): integer loops, same problem. SPEC (1988): real workloads,
standardized methodology. MLPerf (2018): ML-specific, multi-objective.''
Then: ``Each generation solved the gaming problem of its predecessor.''

% -- WARN: Students think gaming is a past problem
Common error: ``Modern benchmarks can't be gamed.'' Correct framing: gaming
evolves with benchmarks. MLPerf addresses known patterns but new ones emerge.

% -- FLEX: [OPTIONAL] Historical context is enriching but not essential
[OPTIONAL] The three cumulative lessons are the key takeaway.
IF AHEAD: ``What gaming patterns does MLPerf still not address?''
IF SHORT: Focus on the three cumulative lessons at the bottom.
}

% --- Layout: FULL-WIDTH IMAGE ---
\centering
\safeimg[width=0.92\textwidth,height=4.8cm,keepaspectratio]{benchmark-evolution-timeline.pdf}

\vspace{0.1cm}
\mlsysinsight{Pattern:}{Each generation of benchmarks solved the gaming problem of its predecessor.}

\end{frame}

% --- ACTIVE LEARNING 1: Predict ---
\begin{frame}{Predict: What Can Go Wrong?}
\note{
% -- LINK: History showed benchmarks get gamed; now students identify failure modes
The timeline showed benchmarks are imperfect. This prediction exercise asks
students to identify production failure modes before seeing the framework.

% -- NARRATE: Frame the scenario, then pause
``A team reports 2x faster with no accuracy loss. List 3 things that could
still go wrong in production.'' Wait 60 seconds. Turn to a neighbor.
Do NOT reveal the three-dimensional framework yet.

% -- ENGAGE: Prediction primes the next slide
Good answers: data drift, preprocessing not measured, batch size unrealistic,
tail latency not reported, thermal throttling, concurrent load effects.

% -- WARN: Students focus on model accuracy, ignore system issues
Common error: all three failures are accuracy-related. Correct framing:
system failures (latency, throughput under load) are equally common.

% -- FLEX: [CORE] Active learning moment 1 of 3
[CORE] Creates need-to-know for the three-dimensional framework.
IF SHORT: Reduce to 30 seconds, share with neighbor.
}

\centering
\vspace{0.8cm}
{\Large\bfseries Think--Write--Share}

\vspace{0.5cm}
{\large A team reports their compressed model is ``2$\times$ faster\\
with no accuracy loss'' on their benchmark.\\[0.2cm]
\alert{List 3 things that could still go wrong in production.}}

\vspace{0.5cm}
{\normalsize Write down 3 failure modes. \textcolor{midgray}{(60 seconds)}}

\vspace{0.3cm}
{\small\textcolor{lightgray}{Turn to a neighbor and compare your answers.}}

\end{frame}

% =============================================================================
\section{Three Dimensions}
% =============================================================================

\begin{frame}{The Three-Dimensional Framework}
\note{
% -- LINK: Students predicted failures; now categorize them into three dimensions
The prediction exercise surfaced failure modes. This framework organizes them:
system, model, and data are orthogonal dimensions.

% -- NARRATE: Walk through the diagram, then connect to predictions
``System benchmarks test hardware and infrastructure. Model benchmarks test
accuracy and robustness. Data benchmarks test distribution and scale.''
Then: ``Which dimension did your prediction-exercise failures fall into?''

% -- ENGAGE: Categorize prediction answers into dimensions
Ask: ``Which dimension did your failures belong to?'' Students realize their
answers were clustered --- the framework reveals blind spots.

% -- WARN: Students think one dimension suffices
Common error: ``I tested accuracy, so I'm done.'' Correct framing: each
dimension reveals failures invisible to the others. System-only misses
compression damage. Model-only misses hardware bottlenecks.

% -- FLEX: [CORE] The framework organizes the rest of the chapter
[CORE] Every subsequent section maps to one or more dimensions.
IF AHEAD: ``Which dimension is hardest to benchmark?''
IF SHORT: Show diagram, state the blind spots.
}

% --- Layout: FULL-WIDTH IMAGE ---
\centering
\safeimg[width=0.92\textwidth,height=4.8cm,keepaspectratio]{three-dimensional-framework.pdf}

\vspace{0.1cm}
\mlsysalert{Key:}{Each dimension reveals failures invisible to the others. Passing all three provides deployment confidence.}

\end{frame}

\begin{frame}{Anatomy of a Benchmark}
\note{
% -- LINK: Three dimensions showed what to measure; anatomy shows how benchmarks work
The framework showed what to test. This slide shows the four layers that make
up any benchmark, like a protocol stack.

% -- NARRATE: Walk through four layers
``Workload (what you run) -> Metrics (what you measure) -> Methodology (how
you measure) -> Reporting (how you present). A benchmark is only as strong
as its weakest layer. A flaw at ANY layer invalidates results.''

% -- WARN: Students trust benchmark results without examining methodology
Common error: accepting numbers without asking how they were obtained.
Correct framing: always check: what workload? what metric? what methodology?

% -- FLEX: [OPTIONAL] Conceptual framework, not a key calculation
[OPTIONAL] The ``weakest layer'' insight is the takeaway.
IF AHEAD: ``Which layer is most commonly flawed?''
IF SHORT: State the four layers and the weakness callout.
}

% --- Layout: FULL-WIDTH IMAGE ---
\centering
\safeimg[width=0.88\textwidth,height=4.8cm,keepaspectratio]{benchmark-anatomy-pipeline.pdf}

\vspace{0.05cm}
{\small\textbf{Four layers}: Workload $\to$ Metrics $\to$ Methodology $\to$ Reporting. A flaw at \emph{any} layer invalidates the results.}

\end{frame}

% =============================================================================
\section{Granularity}
% =============================================================================

\begin{frame}{Benchmarking Granularity}
\note{
% -- LINK: Anatomy showed benchmark layers; granularity shows scope choices
Benchmark anatomy was about structure. Granularity is about scope: how much
of the system do you measure?

% -- NARRATE: Walk through three levels in the diagram
``Micro: individual kernels. Pinpoints bottlenecks but misses interactions.
Macro: model or subsystem. End-to-end: full pipeline including preprocessing.''
Then: ``A 3x kernel speedup may yield 0x system speedup if the data pipeline
cannot keep pace.'' Design rule: match validation to optimization granularity.

% -- ENGAGE: Which granularity for each Ch10 technique?
Ask: ``Which granularity validates kernel fusion? Pruning? Data curation?''
Expected: fusion = micro, pruning = macro, data curation = end-to-end.

% -- WARN: Students benchmark at the wrong granularity
Common error: micro-benchmarking a system-level optimization.
Correct framing: if you optimized the model, validate end-to-end.

% -- FLEX: [CORE] Granularity determines whether you detect real improvements
[CORE] Mismatched granularity is the most common benchmarking error.
IF AHEAD: ``When should you use all three levels?''
IF SHORT: State the three levels and the design rule.
}

% --- Layout: FULL-WIDTH IMAGE ---
\centering
\safeimg[width=0.92\textwidth,height=4.8cm,keepaspectratio]{benchmark-granularity.pdf}

\vspace{0.1cm}
\mlsysconcept{Design rule:}{Match validation granularity to optimization granularity. Kernel fusion $\to$ micro. Pruning $\to$ macro. Data curation $\to$ end-to-end.}

\end{frame}

\begin{frame}{The Latency Breakdown}
\note{
% -- LINK: Granularity showed levels; this waterfall chart shows why end-to-end matters
Granularity was abstract. This concrete example shows what happens when you
optimize one component without measuring end-to-end.

% -- NARRATE: Walk through two waterfall scenarios
``Before: model inference is 60\% of total latency. After TensorRT: 2.5x model
speedup. But preprocessing is now 43\% of total. System speedup: only 1.3x.''
Then: ``This is Amdahl's Law from Ch11 in action.''

% -- ENGAGE: Where should the team invest next?
Ask: ``Where should the team invest next?'' Expected: preprocessing ---
it is now the dominant component (Amdahl's serial fraction).

% -- WARN: Students celebrate component speedups without checking end-to-end
Common error: ``We got 2.5x model speedup!'' Correct framing: the system
got only 1.3x. The bottleneck shifted to preprocessing.

% -- FLEX: [CORE] Amdahl's Law in action --- connects Ch11 to Ch12
[CORE] This is the most impactful slide for changing benchmarking behavior.
IF AHEAD: ``How do you optimize preprocessing?''
IF SHORT: State the 2.5x to 1.3x contrast.
}

% --- Layout: FULL-WIDTH IMAGE ---
\centering
\safeimg[width=0.92\textwidth,height=4.8cm,keepaspectratio]{latency-breakdown-waterfall.pdf}

\vspace{0.1cm}
\mlsysalert{Amdahl's Law:}{2.5$\times$ model speedup $\to$ only 1.3$\times$ system speedup when preprocessing dominates.}

\end{frame}

% =============================================================================
\section{Measurement Rigor}
% =============================================================================

\begin{frame}{Percentiles, Not Averages}
\note{
% -- LINK: Latency breakdown showed total latency; this shows its distribution matters
The waterfall showed mean latency components. This slide reveals that the
mean hides catastrophic tail behavior.

% -- NARRATE: Walk through the distribution diagram
``Mean: 12 ms. But p99: 72 ms --- 6x the mean! At 10,000 QPS, 100 requests
per second hit the tail. With fan-out across 100 services, 63\% of page loads
hit at least one tail event.'' Production SLOs: defined at p95/p99, never mean.

% -- ENGAGE: Why does fan-out make tails worse?
Ask: ``If one service has 1\% chance of tail latency, what happens with 100
services in a fan-out?'' Expected: 1-(0.99)^100 = 63\% chance of at least one.

% -- WARN: Students report mean latency and think they are done
Common error: ``Average latency is 12 ms, meets our SLO.'' Correct framing:
the p99 at 72 ms violates any reasonable SLO. The mean is misleading.

% -- FLEX: [CORE] Tail latency is the production reality
[CORE] This changes how students report benchmarks permanently.
IF AHEAD: ``How do you reduce tail latency?''
IF SHORT: State the 12 ms mean vs 72 ms p99 contrast.
}

% --- Layout: FULL-WIDTH IMAGE ---
\centering
\safeimg[width=0.92\textwidth,height=4.8cm,keepaspectratio]{latency-distribution-percentiles.pdf}

\vspace{0.1cm}
\mlsysalert{Production rule:}{Report p95 and p99, not the mean. At scale, the tail determines the user experience.}

\end{frame}

\begin{frame}{Quick Check}
\note{
% -- NARRATE: Pause for retrieval before measurement methodology
``Close notes. Why report p95/p99 instead of the mean?'' Wait 30 seconds.
Expected: mean hides tail; at scale, tails determine user experience;
production SLOs are defined at percentiles.

% -- FLEX: [CORE] Distributed retrieval between sections
[CORE] 30-second pause cements percentile thinking.
IF SHORT: Reduce to 15 seconds.
}

\centering
\Large\bfseries Quick check --- no peeking.\\[0.5cm]
\normalsize Why should you report p95 and p99 latency instead of the mean?\\[0.3cm]
{\small 30 seconds --- then we continue.}
\end{frame}



\begin{frame}{Measurement Methodology}
\note{
% -- LINK: Percentiles showed what to report; methodology shows how to measure
Percentiles defined the right metric. This slide defines the right protocol
for collecting those measurements.

% -- NARRATE: Walk through the three-phase protocol
``Warm-up: discard first 10--50 iterations (JIT, cache cold start). Measurement
window: collect data under steady state. Cool-down: discard.'' Then the
statistical requirements: ``Confidence intervals, multiple runs (>=5),
environment control.'' Anti-pattern: ``Cherry-picking the best of 3 runs is
marketing, not benchmarking.''

% -- WARN: Students skip warm-up and report best-of-3
Common error: the first iteration is 10x slower due to JIT compilation.
Including it skews the mean. Best-of-3 hides variance.
Correct framing: discard warm-up, report median with 95\% CI from >=5 runs.

% -- FLEX: [CORE] Measurement methodology separates engineering from guesswork
[CORE] Students must follow this protocol in their own benchmarks.
IF AHEAD: ``How many runs are needed for 95\% CI on p99?''
IF SHORT: State warm-up rule and the cherry-picking anti-pattern.
}

% --- Layout: FULL-WIDTH IMAGE ---
\centering
\safeimg[width=0.92\textwidth,height=4.8cm,keepaspectratio]{measurement-methodology.pdf}

\vspace{0.05cm}
{\small\textbf{Never measure first 10--50 iterations.} Report median with 95\% CI from $\geq$5 independent runs.}

\end{frame}

% --- ACTIVE LEARNING 2: Exercise ---
\begin{frame}{Your Turn: The Statistical Confidence Trap}
\note{
% -- LINK: Methodology showed CI is required; this exercise shows why sample size matters
The measurement slide required 95\% confidence intervals. This exercise
demonstrates that small test sets cannot detect small accuracy drops.

% -- NARRATE: Present problem, then reveal after pause
``95\% accuracy model, compressed to 94\%, tested on 1,000 images. Can you
detect the regression?'' Wait 90 seconds. Solution: ``sigma = 6.9. CI = [36,
64]. Both 50 and 60 errors fall inside. 1,000 samples cannot detect 1\%.
Need approximately 10,000 samples.''

% -- ENGAGE: Peer instruction protocol
Present (30s). Individual work (60s). If 30--70\% correct: pair discussion
(90s), re-vote. The key: sample size determines detection power.

% -- WARN: Students trust small test sets for precision claims
Common error: ``94\% on 1,000 images proves the model is worse.'' Correct
framing: both 94\% and 95\% fall within the same confidence interval at N=1000.

% -- FLEX: [CORE] Active learning moment 2 of 3
[CORE] This calculation permanently changes how students evaluate accuracy.
IF AHEAD: ``What if the drop is 5\% instead of 1\%?''
IF SHORT: Show solution immediately.
}

\footnotesize
\begin{columns}[T]
  \begin{column}{0.55\textwidth}
    {\small\bfseries Can You Detect the Regression?}

    \vspace{0.1cm}
    An image classifier has \textbf{95\% accuracy}. A compressed version scores \textbf{94\%} on a \textbf{1,000-image} test set.

    \vspace{0.1cm}
    \textbf{Calculate:}
    \begin{itemize}\setlength\itemsep{0pt}
      \item Expected errors at 95\%: $1000 \times 0.05 = ?$
      \item Std dev: $\sqrt{N \cdot p \cdot (1-p)} = ?$
      \item 95\% CI: $50 \pm 1.96 \times \sigma = ?$
    \end{itemize}

    \vspace{0.05cm}
    Is 60 errors (94\%) outside the CI?

    {\textcolor{midgray}{(90 seconds --- compare with a neighbor)}}
  \end{column}
  \begin{column}{0.42\textwidth}
    \pause
    \begin{mlsyscard}{errorstroke}
    \textbf{Solution:}\\[0.05cm]
    Expected: 50 errors\\
    $\sigma = \sqrt{1000 \times 0.05 \times 0.95} \approx \mathbf{6.9}$\\
    95\% CI: $[36, 64]$\\[0.05cm]
    Both 50 and 60 fall inside!\\[0.05cm]
    \textcolor{errorstroke}{\textbf{1,000 samples cannot detect 1\% drop.}}\\
    Need $\approx$\textbf{10,000} samples.
    \end{mlsyscard}
  \end{column}
\end{columns}

\end{frame}

% =============================================================================
\section{MLPerf}
% =============================================================================

\mlsysfocus{MLPerf}{The industry standard for ML benchmarking:\\[0.2cm]
Representative workloads $+$ Multi-objective metrics $+$ Integrated system measurement}

\begin{frame}{MLPerf Inference Scenarios}
\note{
% -- LINK: Statistical rigor showed how to measure; MLPerf shows the industry standard
Students learned measurement methodology. MLPerf applies it at industry
scale with standardized workloads and scenarios.

% -- NARRATE: Walk through the 2x2 quadrant
``Latency constraint (strict vs relaxed) x input concurrency (single vs
multiple). SingleStream = mobile. Server = cloud API. MultiStream = AV.
Offline = batch scoring.'' Then: ``The scenario determines the metric:
latency for SingleStream, QPS for Server, throughput for Offline.''

% -- ENGAGE: Which scenario for your project?
Ask: ``Which scenario does your ML project fall into?'' Cold-call 2 students.

% -- WARN: Students apply the wrong scenario to their workload
Common error: benchmarking a real-time system with Offline throughput.
Correct framing: the scenario must match the deployment pattern.

% -- FLEX: [CORE] MLPerf scenarios are the industry standard
[CORE] Students will encounter MLPerf in every hardware evaluation.
IF AHEAD: ``Why four scenarios and not just two?''
IF SHORT: Name the four scenarios and their metrics.
}

% --- Layout: FULL-WIDTH IMAGE ---
\centering
\safeimg[width=0.88\textwidth,height=4.8cm,keepaspectratio]{mlperf-scenarios-quadrant.pdf}

\vspace{0.05cm}
{\small The \textbf{scenario} determines the metric: latency for SingleStream, QPS for Server, throughput for Offline.}

\end{frame}

\begin{frame}{MLPerf Training vs.\ Inference}
\note{
% -- LINK: Scenarios showed inference modes; now contrast training vs inference
MLPerf Inference scenarios were just introduced. This slide contrasts training
and inference benchmarking --- fundamentally different optimizations.

% -- NARRATE: Walk through the comparison table row by row
``Training: time-to-quality (reach target accuracy fast). Inference: QPS under
latency SLO. Training: max hardware utilization. Inference: SLO compliance.
Training: large batches. Inference: small--medium.'' Then the cards:
``Training is compute-dominated. Inference is pipeline-dominated.''

% -- WARN: Students apply training benchmarks to inference and vice versa
Common error: reporting training throughput as inference capability.
Correct framing: training optimizes total time; inference optimizes per-query latency.

% -- FLEX: [CORE] Training/inference distinction prevents metric confusion
[CORE] Students must know which benchmark applies to their use case.
IF AHEAD: ``Can the same hardware be optimal for both?''
IF SHORT: Cover primary metric and key bottleneck rows only.
}

\footnotesize
\renewcommand{\arraystretch}{1.05}
\begin{tabular}{@{}lll@{}}
  \toprule
  & \textbf{MLPerf Training} & \textbf{MLPerf Inference} \\
  \midrule
  \textbf{Primary metric}  & Time-to-quality             & QPS under latency SLO \\
  \textbf{Goal}             & Reach target accuracy fast  & Serve queries within SLA \\
  \textbf{Optimization}     & Max hardware utilization    & SLO compliance \\
  \textbf{Batch size}       & Large (throughput)          & Small--medium (latency) \\
  \textbf{Precision}        & Mixed (FP16/BF16 + FP32)   & INT8, FP16, even INT4 \\
  \textbf{Key bottleneck}   & Gradient sync, data loading & Tail latency, preprocessing \\
  \bottomrule
\end{tabular}

\vspace{0.1cm}
\begin{columns}[T]
  \begin{column}{0.48\textwidth}
    \begin{mlsyscard}{computestroke}
    {\footnotesize \textbf{Training}: $T = O/(R_{\text{peak}} \cdot \eta)$ --- compute-dominated}
    \end{mlsyscard}
  \end{column}
  \begin{column}{0.48\textwidth}
    \begin{mlsyscard}{datastroke}
    {\footnotesize \textbf{Inference}: $L = L_{\text{pre}} + L_{\text{model}} + L_{\text{post}}$ --- pipeline-dominated}
    \end{mlsyscard}
  \end{column}
\end{columns}

\end{frame}

\begin{frame}{MLPerf Suite Variants}
\note{
% -- LINK: Training vs inference showed two modes; variants cover all deployments
Training and inference are two of four MLPerf variants, each targeting a
different deployment context.

% -- NARRATE: Walk through the four-variant table
``Training for datacenter. Inference for server/edge. Tiny for MCU/IoT (under
1 mW, under 1 MB). Power as cross-cutting concern.'' Then the crimson card:
``Representative workloads, full-system measurement, open submission.''

% -- WARN: Students think one benchmark suite covers everything
Common error: using MLPerf Inference to evaluate a TinyML deployment.
Correct framing: each variant measures what matters for its deployment target.

% -- FLEX: [OPTIONAL] The variant table is reference material
[OPTIONAL] The design principles are the key takeaway.
IF AHEAD: ``How often does MLPerf update its benchmark suite?''
IF SHORT: Name the four variants, skip the details.
}

\scriptsize
\renewcommand{\arraystretch}{1.1}
\begin{tabular}{@{}lllp{3.5cm}@{}}
  \toprule
  \textbf{Variant} & \textbf{Domain} & \textbf{Constraints} & \textbf{Metrics} \\
  \midrule
  \textbf{Training}  & Datacenter   & Multi-node, high BW       & Time-to-quality, samples/s \\
  \textbf{Inference}  & Server/Edge  & Latency SLAs, throughput  & QPS, latency percentiles \\
  \textbf{Tiny}       & MCU / IoT    & $<$1 mW, $<$1 MB memory  & Latency, energy/inference \\
  \textbf{Power}      & Cross-cutting & Energy budgets, thermal  & Performance/Watt \\
  \bottomrule
\end{tabular}

\vspace{0.15cm}
\footnotesize
\begin{mlsyscard}{crimson}
\textbf{MLPerf design principles}: Representative workloads, full-system measurement, open submission. Vendors must report end-to-end system performance on standardized tasks.
\end{mlsyscard}

\mlsyscite{Reddi et al., MLPerf Inference, 2020; Mattson et al., MLPerf Training, 2020}

\end{frame}

% =============================================================================
\section{System Benchmarks}
% =============================================================================

\begin{frame}{The Fallacy of Peak Performance}
\note{
% -- LINK: MLPerf uses real workloads; this shows why peak specs are misleading
MLPerf requires real workload measurement. This slide quantifies why:
peak performance systematically overpredicts real performance.

% -- NARRATE: Walk through Patterson quote and utilization table
``Peak performance is the performance the manufacturer guarantees you will not
exceed.'' Table: ``A100 FP16: 312 TF peak, 90--155 TF achieved. BERT batch=1:
33\% utilization. DLRM: 40\%.'' Why: ``Memory stalls, pipeline bubbles,
kernel launch overhead, thermal throttling.'' Diagnostic: use the Roofline.

% -- ENGAGE: Calculate the gap for a specific workload
Ask: ``If A100 achieves 33\% utilization on BERT batch=1, what is the actual
TFLOPS?'' Expected: 312 x 0.33 = ~103 TFLOPS.

% -- WARN: Students trust vendor spec sheets
Common error: comparing GPUs by peak TFLOPS in vendor materials.
Correct framing: utilization varies 2--30x depending on workload AI.

% -- FLEX: [CORE] Peak vs achieved is the chapter's central quantitative lesson
[CORE] Students must always ask ``at what utilization?''
IF AHEAD: ``How do vendors game peak numbers?''
IF SHORT: State the Patterson quote and the 33\% BERT number.
}

\footnotesize
\begin{columns}[T]
  \begin{column}{0.55\textwidth}
    \begin{mlsyscard}{errorstroke}
    \textbf{Patterson:} ``Peak performance is the performance the manufacturer \emph{guarantees} you will not exceed.''
    \end{mlsyscard}

    \vspace{0.1cm}
    \renewcommand{\arraystretch}{1.05}
    \begin{tabular}{@{}lrr@{}}
      \toprule
      \textbf{Workload} & \textbf{Peak} & \textbf{Achieved} \\
      \midrule
      A100 FP16 (Tensor)    & 312 TF   & 90--155 TF \\
      ResNet-50 (batch 256) & ---      & 85--90\% util. \\
      BERT (batch 1)        & ---      & $\sim$33\% util. \\
      DLRM (rec.)           & ---      & $\sim$40\% util. \\
      \bottomrule
    \end{tabular}

    \vspace{0.05cm}
    The gap is 2--3$\times$ even in optimally tuned systems.
  \end{column}
  \begin{column}{0.42\textwidth}
    \textbf{Why the gap?}
    \begin{itemize}\setlength\itemsep{0pt}
      \item \textcolor{datastroke}{Memory stalls}: data starves the ALU
      \item \textcolor{computestroke}{Pipeline bubbles}: sequential deps
      \item \textcolor{routingstroke}{Kernel launch}: dispatch overhead
      \item \textcolor{errorstroke}{Thermal throttling}: sustained load
    \end{itemize}

    \vspace{0.05cm}
    \mlsysconcept{Diagnostic:}{Use the \textbf{Roofline Model} to identify compute-bound vs.\ memory-bound.}
  \end{column}
\end{columns}

\end{frame}

\begin{frame}{Roofline: ResNet vs.\ BERT}
\note{
% -- LINK: Peak performance showed the gap exists; Roofline diagnoses it
The fallacy slide showed utilization varies wildly. The Roofline from Ch11
explains why: arithmetic intensity determines the ceiling.

% -- NARRATE: Walk through the table, then the batching insight
``ResNet batch=256: AI=300, compute-bound, 85--90\% util. BERT batch=1:
AI=50, memory-bound, 33\% util.'' Then the key insight: ``Batching transforms
BERT from memory-bound to compute-bound: batch=32 pushes AI to 1600,
crossing the ridge. 85\% utilization.'' Tradeoff: higher throughput but higher
latency (must accumulate requests).

% -- ENGAGE: Connect to MLPerf scenarios
Ask: ``Which MLPerf scenario captures this throughput-latency tradeoff?''
Expected: Server scenario --- measures QPS under a latency SLO.

% -- WARN: Students forget batching changes the regime
Common error: ``BERT is always memory-bound.'' Correct framing: it depends
on batch size. Batch=1 is deeply memory-bound; batch=32 is compute-bound.

% -- FLEX: [CORE] This connects Ch11's Roofline to Ch12's benchmarking
[CORE] The bridge between theory (Roofline) and practice (MLPerf).
IF AHEAD: ``What is the optimal batch size for latency vs throughput?''
IF SHORT: State the batch=1 vs batch=32 contrast.
}

\footnotesize
\textbf{A100 GPU}: 312 TFLOPS FP16, 2.0 TB/s bandwidth, Ridge $\approx$ 153 FLOPs/byte

\vspace{0.05cm}
\renewcommand{\arraystretch}{1.05}
\begin{tabular}{@{}lrrrl@{}}
  \toprule
  \textbf{Workload} & \textbf{AI (FLOPs/B)} & \textbf{Regime} & \textbf{Util.} & \textbf{Diagnosis} \\
  \midrule
  ResNet-50 (batch 256) & $\sim$300 & Compute-bound & 85--90\% & $\checkmark$ Efficient \\
  BERT-Base (batch 1)   & $\sim$50  & Memory-bound  & $\sim$33\% & $\times$ Underutilized \\
  BERT-Base (batch 32)  & $\sim$1600 & Compute-bound & $\sim$85\% & $\checkmark$ Batching helps \\
  \bottomrule
\end{tabular}

\vspace{0.1cm}
\begin{mlsyscard}{datastroke}
\textbf{Batching transforms the regime}: BERT at batch 1 is memory-bound (33\% util.). At batch 32, same weights serve 32$\times$ more compute $\to$ 85\% utilization.
\textbf{Tradeoff}: Higher throughput but higher latency (must accumulate requests).
\end{mlsyscard}

\end{frame}

\begin{frame}{Decoding Vendor Claims}
\note{
% -- LINK: Roofline diagnosed real performance; this slide decodes marketing language
The Roofline showed what performance is achievable. This translation guide
helps students decode what vendors actually mean.

% -- NARRATE: Walk through each claim in the table
``Up to 10,000 images/s: peak throughput at max batch, INT8, no preprocessing.
Sub-millisecond latency: accelerator compute only. 1000 TOPS: at what
precision?'' Then: ``Always ask: what is measured, what is excluded, what
conditions?''

% -- WARN: Students accept vendor claims at face value
Common error: using marketing numbers in system design calculations.
Correct framing: every claim has hidden assumptions. Decode them.

% -- FLEX: [OPTIONAL] Practical but not conceptually deep
[OPTIONAL] The ``always ask'' checklist is the key takeaway.
IF AHEAD: ``How would you independently verify a vendor claim?''
IF SHORT: Read the ``always ask'' checklist.
}

\scriptsize
\renewcommand{\arraystretch}{1.1}
\begin{tabular}{@{}lp{6.5cm}@{}}
  \toprule
  \textbf{Vendor Claim} & \textbf{What It Often Means} \\
  \midrule
  ``Up to 10,000 images/sec'' & Peak throughput at max batch, INT8, no preprocessing \\
  ``Sub-millisecond latency'' & Accelerator compute only, excluding data transfer \\
  ``5$\times$ more efficient'' & Per-operation efficiency, not total system efficiency \\
  ``1000 TOPS'' & At what precision? INT4 TOPS are 4$\times$ INT8 on same HW \\
  ``2$\times$ faster'' & On which workload? What batch size? What precision? \\
  \bottomrule
\end{tabular}

\vspace{0.1cm}
\footnotesize
\textbf{Always ask:} What is measured? (Model-only or end-to-end?) What is excluded? (Memory transfer, throttling?) What conditions? (Batch size, precision, thermal state?)

\end{frame}

% --- ACTIVE LEARNING 3: Discussion ---
\begin{frame}{Discussion: Diagnose This Benchmark}
\note{
% -- LINK: Students learned the Roofline and vendor decoding; now apply both
Students know the Roofline (Ch11), peak vs achieved, and vendor decoding.
This exercise integrates all three.

% -- NARRATE: Frame the problem, then facilitate discussion
``BERT inference on H100, 20\% utilization. Is it compute-bound or memory-bound?
What to optimize? Will a faster GPU help?'' 90 seconds pair discussion.
Cold-call 2--3 pairs. Answer: memory-bound (AI < ridge). More TFLOPS won't
help. Increase batch size or quantize to reduce memory traffic.

% -- ENGAGE: Force Roofline-based reasoning
Students must use AI and ridge point to justify their answer. Vague answers
(``it's slow'') are insufficient.

% -- WARN: Students default to ``buy faster hardware''
Common error: ``20\% util means we need a bigger GPU.'' Correct framing:
20\% util means the GPU is starved for data. More compute makes it worse.

% -- FLEX: [CORE] Active learning moment 3 of 3 --- synthesis
[CORE] Integrates Roofline, benchmarking, and optimization.
IF AHEAD: ``What batch size would push utilization to 80\%?''
IF SHORT: Reduce to 60 seconds, cold-call one pair.
}

\centering
\vspace{0.6cm}
{\large\bfseries Turn and Talk \textcolor{midgray}{(90 seconds)}}

\vspace{0.5cm}
{\large A team benchmarks BERT inference on an H100.\\
They achieve only \textbf{20\% GPU utilization}.\\[0.3cm]
\alert{Using the Roofline Model (Ch 11):}}

\vspace{0.4cm}
\begin{columns}[c]
  \begin{column}{0.30\textwidth}\centering\small Is it compute-bound\\or memory-bound?\end{column}
  \begin{column}{0.30\textwidth}\centering\small What should they\\optimize first?\end{column}
  \begin{column}{0.30\textwidth}\centering\small Will buying a\\faster GPU help?\end{column}
\end{columns}

\end{frame}

% =============================================================================
\section{Power \& Energy}
% =============================================================================

\begin{frame}{Power Measurement Techniques}
\note{
% -- LINK: System benchmarks covered throughput/latency; energy is the third dimension
Throughput and latency were covered. This adds energy as a first-class
benchmark dimension --- critical for edge and mobile deployment.

% -- NARRATE: Walk through both columns
Left: ``Why energy matters: cloud (cost), edge (15--40W thermal), mobile
(3--5W battery), TinyML (100mW coin-cell).'' INT8 card: ``4x memory reduction
but 10--20x energy reduction --- DRAM to SRAM shift.'' Right: ``System
boundary matters: chip, board, server, rack.'' MLPerf Power: 11 load levels.

% -- ENGAGE: Why is the system boundary important?
Ask: ``If you measure only the chip, what do you miss?'' Expected: memory
system, VRM losses, cooling, networking. Chip-only can understate total
energy by 2--5x.

% -- WARN: Students measure energy at the wrong boundary
Common error: reporting chip-only power as system power. Correct framing:
production energy includes everything in the rack.

% -- FLEX: [CORE] Energy benchmarking is mandatory for edge/mobile
[CORE] Students deploying to edge must measure energy, not just latency.
IF AHEAD: ``How does idle power affect total cost of ownership?''
IF SHORT: Cover INT8 energy lever and system boundary table.
}

\small
\begin{columns}[T]
  \begin{column}{0.52\textwidth}
    \textbf{Why energy matters:}
    \begin{itemize}\setlength\itemsep{0pt}
      \item {\footnotesize Cloud: energy cost rivals hardware cost}
      \item {\footnotesize Edge: 15--40 W thermal envelope}
      \item {\footnotesize Mobile: 3--5 W battery constraint}
      \item {\footnotesize TinyML: 100 mW coin-cell budget}
    \end{itemize}

    \vspace{0.1cm}
    \begin{mlsyscard}{datastroke}
    \textbf{INT8 as energy lever}\\
    {\footnotesize 4$\times$ memory reduction but 10--20$\times$ energy reduction: DRAM access ($\sim$20 pJ/bit) $\to$ SRAM + INT8 ALU ($\sim$1 pJ/bit).}
    \end{mlsyscard}
  \end{column}
  \begin{column}{0.45\textwidth}
    \textbf{System boundary matters:}

    \vspace{0.1cm}
    {\footnotesize
    \renewcommand{\arraystretch}{1.15}
    \begin{tabular}{@{}ll@{}}
      \toprule
      \textbf{Boundary} & \textbf{Includes} \\
      \midrule
      Chip only   & Accelerator die \\
      Board       & + memory, VRM \\
      Server      & + CPU, fans, PSU \\
      Rack        & + networking, cooling \\
      \bottomrule
    \end{tabular}
    }

    \vspace{0.15cm}
    {\footnotesize MLPerf Power measures at server boundary across 11 load levels (0\%--100\% in 10\% steps).}

    \vspace{0.1cm}
    {\footnotesize\textcolor{midgray}{Idle power matters: real utilization is 20--50\%.}}
  \end{column}
\end{columns}

\end{frame}

% =============================================================================
\section{Wrap-Up}
% =============================================================================

\begin{frame}{Benchmark Gaming: Anti-Patterns}
\note{
% -- LINK: Measurement rigor was defensive; gaming is the offensive threat
Measurement methodology defended against honest mistakes. Benchmark gaming
is deliberate manipulation --- the threat model.

% -- NARRATE: Walk through anti-patterns in the diagram
``Cherry-picking: best of 3 runs. Unfair comparison: tuned vs default.
Misleading metric: mean hides tail. Narrow workload: one model doesn't
represent production.'' Then: ``MLPerf resists all four by requiring
diverse suites, percentile reporting, and open submission rules.''

% -- WARN: Students may unknowingly commit gaming anti-patterns
Common error: reporting best run because ``it represents the system's
capability.'' Correct framing: the best run represents the best case.
The median of multiple runs represents typical performance.

% -- FLEX: [CORE] Gaming awareness prevents both committing and falling for it
[CORE] Students must recognize these patterns in vendor benchmarks.
IF AHEAD: ``What new gaming patterns could emerge for LLM benchmarks?''
IF SHORT: Name four anti-patterns, state how MLPerf resists them.
}

% --- Layout: FULL-WIDTH IMAGE ---
\centering
\safeimg[width=0.92\textwidth,height=4.8cm,keepaspectratio]{benchmark-gaming-antipatterns.pdf}

\vspace{0.05cm}
{\small MLPerf resists gaming by requiring \textbf{diverse suites}, \textbf{percentile reporting}, and \textbf{open submission rules}.}

\end{frame}

\begin{frame}{Fallacies}
\note{
% -- LINK: Gaming was about manipulation; fallacies are about honest misconceptions
Gaming was deliberate. These fallacies are honest mistakes that even
well-intentioned engineers make.

% -- NARRATE: Walk through three fallacies with numbers
First (most important): ``92\% benchmark accuracy drops to 78--82\% in
production. 15 ms mean becomes 150--200 ms p99 under load.'' Second:
``94\% accuracy but 180 ms p99. SLO requires p99 < 100 ms. Accurate but
unusable.'' Third: ``312 TFLOPS peak, 90--155 sustained. Even optimally tuned.''

% -- WARN: The benchmark-to-production fallacy is the most common
Common error: treating benchmark accuracy as production accuracy.
Correct framing: production includes data drift, load, preprocessing.

% -- FLEX: [CORE] Fallacies synthesize the chapter's lessons
[CORE] Each fallacy references a specific concept from the lecture.
IF AHEAD: ``How do you budget for the benchmark-production gap?''
IF SHORT: Cover the first fallacy only.
}

\small
\textbf{Fallacy:} \textit{Benchmark performance directly translates to production.}\\
{\footnotesize Language model: 92\% benchmark accuracy $\to$ 78--82\% on user-generated text. Inference: 15 ms mean $\to$ 150--200 ms p99 under load (10$\times$ degradation).}

\vspace{0.15cm}
\textbf{Fallacy:} \textit{Single-metric evaluation suffices.}\\
{\footnotesize Rec.\ model: 94\% accuracy but 180 ms p99 latency. SLO requires p99 $<$ 100 ms. Accurate but \emph{unusable}. Throughput-optimized system: 1,200 QPS at 420 W vs.\ 1,000 QPS at 180 W. For edge: 17\% throughput loss enables 2.3$\times$ longer battery life.}

\vspace{0.15cm}
\textbf{Fallacy:} \textit{Peak performance predicts deployment.}\\
{\footnotesize A100: 312 TFLOPS peak, 90--155 TFLOPS sustained (30--50\% MFU). The gap exists even in optimally tuned systems.}

\end{frame}

\begin{frame}{Pitfalls}
\note{
% -- LINK: Fallacies were misconceptions; pitfalls are operational mistakes
Fallacies were wrong beliefs. Pitfalls are wrong practices --- things teams
actually do that waste time and money.

% -- NARRATE: Walk through three pitfalls
First: ``ImageNet saturated at 95\% by 2017. MNIST at 99.8\%. Retire
saturated benchmarks.'' Second: ``Latency 12 ms to 8 ms, climbed 15 positions
on MLPerf --- while degrading calibration and increasing energy 40\%. Won
the leaderboard, destroyed the product.'' Third: ``800 QPS isolated, 400--500
under 90\% utilization. Research ignores availability requirements.''

% -- WARN: Goodhart's Law in action
Common error: optimizing for the leaderboard instead of the product.
Correct framing: the benchmark is a proxy. When it becomes the goal,
it stops being useful.

% -- FLEX: [CORE] Pitfalls are mistakes students will encounter
[CORE] These are real-world errors from industry experience.
IF AHEAD: ``How do you know when a benchmark is saturated?''
IF SHORT: Cover the first two pitfalls.
}

\footnotesize
\textbf{Pitfall:} \textit{Using outdated benchmarks that no longer discriminate.}\\
ImageNet: 28.2\% error (2010) $\to$ 3.57\% (2015). Competition ended 2017 when 29/38 teams exceeded 95\%. MNIST: 99.8\% with simple models. Further optimization has marginal value. \textbf{Retire saturated benchmarks.}

\vspace{0.15cm}
\textbf{Pitfall:} \textit{Optimizing exclusively for benchmark metrics.}\\
A team reduces latency from 12 ms to 8 ms via aggressive quantization, climbing 15 positions on MLPerf---while degrading calibration and increasing energy by 40\%. ``Won the leaderboard but destroyed the product.''

\vspace{0.15cm}
\textbf{Pitfall:} \textit{Applying research benchmarks to evaluate production systems.}\\
Research: 800 QPS isolated. Production: 400--500 QPS under 90\% utilization (38\% degradation). Research reports model inference (5--10 ms); production end-to-end is 50--100 ms. Production requires 99.9\% availability (43 min downtime/month)---research benchmarks ignore this entirely.

\end{frame}


\begin{frame}{Muddiest Point}
\note{
% -- NARRATE: Collect muddiest points
``Write the single most confusing concept. One sentence.'' Use to open
the next lecture (Model Serving) with targeted clarification.

% -- FLEX: [CORE] Do not skip
[CORE] 60 seconds, irreplaceable feedback.
IF SHORT: Reduce to 30 seconds.
}

\centering
\Large\bfseries One-minute paper\\[0.5cm]
\normalsize Write down the \textbf{muddiest point} from today's lecture.\\[0.2cm]
{\small What concept was most confusing or unclear?}\\[0.5cm]
{\footnotesize\textcolor{midgray}{(Hand in on your way out --- or submit digitally)}}
\end{frame}

% --- RETRIEVAL PRACTICE ---
\begin{frame}{What Were the Key Ideas?}
\note{
% -- NARRATE: Retrieval practice before takeaways
``Close notes. Write 4 most important concepts. 90 seconds.'' Walk around.
Do NOT show Key Takeaways until time is up.

% -- FLEX: [CORE] Highest-impact learning activity
[CORE] Even 60 seconds doubles retention.
IF SHORT: Reduce to 60 seconds.
}

\centering
\vspace{1.5cm}
{\Large\bfseries Close your notes.}

\vspace{0.8cm}
{\large Write down the \textbf{4 most important concepts} from today.}

\vspace{0.8cm}
{\normalsize\textcolor{midgray}{90 seconds --- no peeking.}}

\end{frame}

\begin{frame}{Key Takeaways}
\note{
% -- NARRATE: Reveal after retrieval, anchor with numbers
``2--10x gap. 1\% detection needs 10K samples. p99 = 6x mean. 2.5x model =
1.3x system. 10--20x energy from INT8. Three dimensions. Goodhart's Law.''

% -- FLEX: [CORE] Definitive summary
[CORE] Students compare retrieval answers against these bullets.
IF SHORT: Read bold terms only.
}

\scriptsize
\begin{itemize}\setlength\itemsep{0pt}
  \item \textbf{Three-dimensional framework}: System, model, and data benchmarks each test different failure modes. Full validation requires all three.
  \item \textbf{Benchmarks are proxies}: Production depends on data distribution, load patterns, and SLA constraints absent from controlled evaluation.
  \item \textbf{Granularity determines diagnostic power}: Micro pinpoints bottlenecks; macro reveals subsystem issues; end-to-end validates deployment.
  \item \textbf{The tail determines user experience}: p99 can be 6$\times$ the mean. Report percentiles, not averages.
  \item \textbf{Amdahl's Law sets the ceiling}: 2.5$\times$ model speedup $\to$ 1.3$\times$ system speedup when preprocessing dominates.
  \item \textbf{Statistical rigor}: 1,000 samples cannot detect a 1\% accuracy drop. Warm-up, CI, and multiple runs are mandatory.
  \item \textbf{INT8 is an energy lever}: 4$\times$ memory reduction but 10--20$\times$ energy reduction (DRAM $\to$ SRAM + INT8).
\end{itemize}

\end{frame}

\begin{frame}{References}
\note{
% -- NARRATE: Point to canonical papers
``Reddi and Mattson for MLPerf, Coleman for DAWNBench, Williams for
Roofline, Henderson for reproducibility, Sculley for technical debt.''

% -- FLEX: [OPTIONAL] Reference slide
[OPTIONAL] Not a teaching moment.
IF SHORT: Say ``references on this slide'' and advance.
}

\small
\mlsysref{Reddi+20}{V.\ J.\ Reddi et al. ``MLPerf Inference Benchmark.'' ISCA, 2020.}
\mlsysref{Mattson+20}{P.\ Mattson et al. ``MLPerf Training Benchmark.'' MLSys, 2020.}
\mlsysref{Coleman+17}{C.\ Coleman et al. ``DAWNBench: End-to-End Deep Learning Benchmark.'' 2017.}
\mlsysref{Williams+09}{S.\ Williams et al. ``Roofline: An Insightful Visual Performance Model.'' Comm.\ ACM, 2009.}
\mlsysref{Henderson+18}{P.\ Henderson et al. ``Deep Reinforcement Learning that Matters.'' AAAI, 2018.}
\mlsysref{Sculley+15}{D.\ Sculley et al. ``Hidden Technical Debt in ML Systems.'' NeurIPS, 2015.}

\end{frame}

\begin{frame}{Next Lecture: Model Serving}
\note{
% -- LINK: This chapter validated in the lab; next meets real traffic
We validated optimizations under controlled conditions. Model serving is
where benchmarks become SLAs and laboratory numbers meet production chaos.

% -- NARRATE: Build anticipation for the next chapter
``We validated in the lab. Now those optimizations meet real traffic: bursts,
variable inputs, concurrent users. Batching trades latency for throughput.
Load balancing minimizes tail latency. SLA compliance is where benchmarks
become contracts.''

% -- FLEX: [CORE] Forward hook motivates next attendance
[CORE] End on: ``From lab to live.''
IF SHORT: State the forward hook and dismiss.
}

\footnotesize
\begin{columns}[T]
  \begin{column}{0.30\textwidth}
    \centering
    \begin{mlsyscard}{computestroke}
    {\small\bfseries Batching}\\[0.05cm]
    {\scriptsize Dynamic batching trades latency for throughput}
    \end{mlsyscard}
  \end{column}
  \begin{column}{0.30\textwidth}
    \centering
    \begin{mlsyscard}{routingstroke}
    {\small\bfseries Load Balancing}\\[0.05cm]
    {\scriptsize Route requests to minimize tail latency}
    \end{mlsyscard}
  \end{column}
  \begin{column}{0.30\textwidth}
    \centering
    \begin{mlsyscard}{errorstroke}
    {\small\bfseries SLA Compliance}\\[0.05cm]
    {\scriptsize Where benchmarks become contracts}
    \end{mlsyscard}
  \end{column}
\end{columns}

\vspace{0.2cm}
\centering
{\normalsize From lab to live: where benchmark numbers become SLAs.}\\[0.05cm]
{\footnotesize\textbf{Production reality includes traffic bursts, data drift, and cascading failures no static benchmark captures.}}

\end{frame}


% =============================================================================
% BACKUP SLIDES
% =============================================================================
\appendix

\begin{frame}{Backup: Confidence Interval Calculation}
\note{
% -- NARRATE: Extended CI calculation for additional practice
For p=0.95 on N=1,000: sigma = sqrt(47.5) = 6.9 errors. 95\% CI: [36, 64].
Cannot distinguish 95\% from 94\%. Need N = 10,000 for 1\% precision.

% -- FLEX: [OPTIONAL] Deploy during Q&A or review
[OPTIONAL] Backup slide.
}
\small
Use if students need more practice with statistical rigor.\\small For accuracy $p = 0.95$ on $N = 1{,}000$ samples:\\ $\sigma = \sqrt{N \cdot p \cdot (1-p)} = \sqrt{47.5} \approx 6.9$ errors.\\ 95\% CI: $50 \pm 1.96 \times 6.9 = [36, 64]$ errors.\\ Cannot distinguish 95\% from 94\% (both inside CI).\\ Need $N \approx 10{,}000$ to detect 1\% differences.
\end{frame}

\begin{frame}{Backup: MLPerf Scenario Selection Guide}
\note{
% -- NARRATE: Quick-reference for scenario selection
SingleStream: mobile, metric = latency. Server: cloud API, metric = QPS.
MultiStream: AV sensors, metric = streams at latency. Offline: batch,
metric = samples/second.

% -- FLEX: [OPTIONAL] Deploy when students ask which scenario applies
[OPTIONAL] Backup reference.
}
\small
Use if students ask which MLPerf scenario applies to their project.\\small \textbf{SingleStream}: One request at a time (mobile, edge). Metric: latency.\\ \textbf{Server}: Poisson arrivals (cloud API). Metric: QPS under latency SLO.\\ \textbf{MultiStream}: Synchronized sensors (AV). Metric: streams at latency.\\ \textbf{Offline}: All data available (batch). Metric: samples/second.
\end{frame}

\end{document}
